% ECONOMICS_PROBLEM: {"id": "F33-355", "title": "Reserve-currency persistence", "jel_code": "F33", "source_id": "OP-0334"}
% !TeX program = xelatex
\documentclass[11pt,a4paper]{article}
\usepackage[margin=23mm]{geometry}
\usepackage{fontspec}
\setmainfont{Latin Modern Roman}
\usepackage{amsmath,amssymb,mathtools}
\usepackage{xcolor,enumitem,booktabs,longtable,array,xurl,hyperref}
\definecolor{muted}{HTML}{53636C}
\hypersetup{colorlinks=true,urlcolor=blue,linkcolor=blue}
\setlength{\parindent}{0pt}
\setlength{\parskip}{5pt}
\newcommand{\R}{\mathbb R}
\newcommand{\E}{\mathbb E}
\newcommand{\N}{\mathbb N}
\newcommand{\ind}{\mathbf 1}
\newcommand{\Var}{\operatorname{Var}}
\newcommand{\Cov}{\operatorname{Cov}}
\newcommand{\supp}{\operatorname{supp}}
\DeclareMathOperator*{\argmin}{arg\,min}
\DeclareMathOperator*{\argmax}{arg\,max}
\newcommand{\entrymeta}[3]{{\sffamily\small\color{muted}\textbf{\texttt{#1}}\quad|\quad #2\par #3\par}}
\newcommand{\Problemref}[1]{\texttt{#1}}
\newcommand{\JELref}[2]{\texttt{#2} (JEL \texttt{#1})}
\begin{document}
\section*{Reserve-currency persistence}
\textbf{Problem ID:} \texttt{F33-355}\quad \textbf{JEL:} \texttt{F33}\par
% BEGIN ECONOMICS STATEMENT
\label{problem:OP-0334}
{\sffamily\small\color{muted}Original subject group: International finance and exchange rates\par}
\label{definitions:problem:30007616}
\textbf{Audit classification:} Published research agenda. The agenda identifies currency evolution as open. The reserve-share threshold, network payoff and transition experiment are editorial; reserve portfolios and invoicing have distinct targets.
\subsection*{Definitions and precise research target}
Consider central banks choosing reserve portfolios across currencies \(j=1,\ldots,J\). Currency \(j\) offers returns, liquidity, settlement compatibility and exposure to issuer fiscal or sanctions risk. Let \(v_j\) be its equilibrium reserve share, with \(v_j\geq0\) and \(\sum_jv_j=1\). A reserve holder's payoff from currency \(j\) depends on fundamentals \(x_j\) and an explicitly specified network-benefit function \(n_j(v_j)\). Digital payment technologies alter settlement costs and feasible choices; these changes are primitives, not assumed conclusions. The task is to characterize equilibria and transition thresholds at which a persistent shock to \(x\) or payment technology substantially lowers the dollar share. Fix a material-decline threshold \(\eta\in(0,1)\) and horizon \(H\), and study
\[\Pr(v_{\mathrm{USD},t+H}\leq v_{\mathrm{USD},t}-\eta\mid do(x_t=x')).\]
Here USD denotes the dollar and \(x'\) is a fully specified intervention. Identify the role of coordination, switching costs and available substitutes, including multiple equilibria and hysteresis. Discipline network-benefit and risk parameters using reserve-holder behavior rather than assuming dominance is permanent. Distinguish reserve holdings from trade invoicing and private funding roles. Predicting one function's decline does not establish collapse of every international role of the currency.

\textbf{Additional formal definitions and scope (editorial).} A reserve portfolio comprises official claims valued consistently by currency; \(v_j\) uses a declared denominator and valuation-date convention. Give reserve-holder \(i\)'s feasible simplex \(A_i\), switching cost \(C_i(v_i,v_i^{\rm old})\), and objective \(U_i(v_i;x,v)\), where aggregate \(v\) is a weighted sum of choices. A static equilibrium solves optimal portfolio choices and aggregation; a dynamic threshold requires a transition law for fundamentals, learning and adjustments. The event threshold is an absolute percentage-point drop \(\eta\), and is feasible only when \(v_{\rm USD,t}\geq\eta\); a proportional decline is a different target. With multiple equilibria, give a selection mechanism or bounds over all admissible equilibrium transition laws. Network coordination and strategic issuer responses are game-theoretic extensions; reserve-share causal estimation can be posed without claiming that canonical open-game results are involved.
\subsection*{Variants and assumption changes}
\begin{enumerate}
\item \textbf{Fixed settlement technology (editorial).} Estimate reserve demand from returns, issuer risk and switching costs with no network payoff.
\item \textbf{Network technology transitions (editorial; strategic extension).} Add a defined network benefit and dynamic selection rule; bound the probability of an absolute reserve-share decline after a specified persistent intervention.
\end{enumerate}
\subsection*{References and source audit}
\begin{enumerate}
\item Official January 9, 2026 web version; locator: Interconnected World / International finance, global-currency paragraph. Broad agenda; exact model editorial. URL: \url{https://www.bankofengland.co.uk/research/bank-of-england-agenda-for-research}.
\end{enumerate}
\begin{enumerate}\raggedright
\item\hypertarget{definitions:source:30007616:1}{} Bank of England. 2026. \emph{Bank of England Agenda for Research, 2025–2028}. Interconnected World / International finance, global-currency paragraph.
\url{https://www.bankofengland.co.uk/research/bank-of-england-agenda-for-research}.
\end{enumerate}

\subsection*{Common conventions: Source mathematical conventions}

These conventions apply to each entry unless explicitly replaced.

\textbf{Models and identification.} A primitive model \(m\) specifies all probability laws, feasible choices, information, equilibrium and measurement rules used by its target. The admissible class \(\mathcal M\) is fixed independently of outcomes. If \(P_O(m)\) is its observable probability law and \(\tau(m)\) the target, its identified set at observable law \(P\) is
\[
 \mathcal I_\tau(P)=\{\tau(m):m\in\mathcal M,\ P_O(m)=P\}.
\]
Point identification means that this set is a singleton. An empty set signals incompatibility of data and assumptions. Coordinate bounds need not describe the joint set. Bounds are sharp only if every retained target is attainable or arbitrarily approachable by compatible models. Finite bounds require explicit restrictions on unobserved quantities; unrestricted confounding, hidden wealth, extrapolation or arbitrary equilibrium selection can yield uninformative sets.

\textbf{Causal interventions.} A potential outcome \(Y_i(p)\) is the outcome under a complete policy regime \(p\), including timing, financing, information and market spillovers. The target population and units are fixed before treatment. With interference, use the complete assignment vector or a justified exposure mapping; individual no-interference notation alone cannot identify an economy-wide intervention. A difference of mean potential outcomes does not require the cross-world joint distribution. Conditioning on post-treatment migration, survival, enrollment or employment changes the estimand and needs separate justification.

\textbf{Statistical statements.} An estimator, confidence set, forecast or test specifies a sampling law, model class and loss/coverage criterion. Uniform \((1-\alpha)\)-coverage means \(\inf_{m\in\mathcal M}\Pr_m\{\tau(m)\in C_n\}\geq1-\alpha\), or an explicitly asymptotic version. The whole parameter space trivially has coverage, so informativeness requires a stated width, risk or power objective. A forecasting improvement is relative to a named benchmark, information set and evaluation population.

\textbf{Equilibrium and optimization.} An equilibrium comprises feasible plans, optimality, beliefs where needed, budget constraints and all market/resource clearing. The primitives or a stated benchmark define each correspondence \(\mathcal E(p;m)\). Nonemptiness is assumed or proved, never inferred just from compact policy sets. A selected-equilibrium target uses a declared selection rule; a robust target quantifies over all permitted equilibria. Use suprema or infima unless attainment is established. An \(\varepsilon\)-optimal policy is feasible and within \(\varepsilon\) of the finite optimum value. Report an empty feasible set separately.

\textbf{Welfare and accounting.} Normative utility, interpersonal normalization, social weights, numeraire, discounting and the use of public receipts are primitives. Behavioral utility need not coincide with the welfare criterion. Household welfare includes ownership income and rents once; adding profits again to compensating variation generally double-counts them. A monetary equivalent is defined by an explicit transfer and price convention, with monotonicity and range conditions. Average effects, mechanism contributions, transfers and welfare losses are different targets. Infinite sums require convergence and dynamic feasible plans require terminal or no-Ponzi conditions.

\textbf{Source versions.} A working-paper date, revision date and journal publication year are distinguished. A DOI for a journal version is not automatically the DOI of an earlier manuscript. Locators refer to the stated version and printed pages where specified. A metadata-only check does not verify a claim about a full-text conclusion. Every entry's source subsection narrows the provenance claim accordingly.
% END ECONOMICS STATEMENT
\end{document}
